\section{Compliance matrix to software product assurance requirements}
{\color{pink} \it
a. The SPAP shall include the compliance matrix to the applicable software
product assurance requirements (e.g. ECSS-Q-ST-80 clauses, as tailored
by a product assurance requirements document), or provide a reference
to it.

b. For each software product assurance requirement, the following
information shall be provided:
1. requirement identifier;
2. compliance
(C = compliant, NC = non–compliant, NA = not applicable);
3. reference to the project documentation covering the requirement
(e.g. section of the software product assurance plan);
4. remarks.
}

The following table shows the compliance matrix to the software product assurance requirements from ECSS tailoring [ECSS-T]. Compliance can be C (compliant), NC (non–compliant), NA (not applicable).

\newcolumntype{L}{>{\hsize=0.4\hsize\raggedright\arraybackslash}X}
\newcolumntype{S}{>{\hsize=0.15\hsize\raggedright\arraybackslash}X}
\newcolumntype{N}{>{\raggedright\arraybackslash}X}

{
    \renewcommand{\arraystretch}{1.4} % introduce padding on top and bottom of cells
    \tiny
    \fontfamily{phv}\selectfont
    
    \begin{xltabular}{\textwidth}{|l|L|l|S|S|}
        % Header on all pages
        \hline
        ID & Requirement & Compliance & Document reference & Remarks \\ \hline\hline
        \endhead

        SPA-01 & 
        {\textbf{Organization and responsibility} \newline
        The Supplier shall: \newline
        • ensure that an organizational structure is defined for software development, that individuals have defined tasks and responsibilities, that their skills are suitable for the activities to be carried out; \newline
        • define and document responsibilities, authority, interfaces of personnel whose work affects software quality; \newline
        • provide resources to perform software product assurance tasks; \newline
        • identify the personnel responsible for software product assurance; \newline
        • grant the software product assurance personnel organisational authority and independence, as well as unimpeded access to higher management as necessary to fulfil their duties.} &
        &
        {PMP, \newline SDP 4.7, \newline SPAP 5.1, 5.2, 5.3} & 
        \\ \hline

        SPA-02 &
        {\textbf{Software product assurance planning and control} \newline
        The Supplier shall develop a software product assurance plan in response to the applicable software product assurance requirements. The software product assurance plan shall be part of the overall product assurance plan.} &
        C &
        SPAP &
        \\ \hline

        SPA-03 &
        {\textbf{Software product assurance reporting} \newline
        The Supplier shall report to their customer on a regular basis on the status of the software product assurance programme implementation, as part of the overall project reporting. \newline
        The software product assurance report shall include: \newline
        • an assessment of the current quality of the product and processes, based on measured properties, with reference to the metrication as defined in the software product assurance plan; \newline
        • product assurance verifications undertaken; \newline
        • problems detected; \newline
        • problems resolved.} &
        C &
        {SPAP 5.4, \newline SPAMR} &
        \\ \hline

        SPA-04 &
        {\textbf{Audits} \newline
        The Supplier shall perform audits on lower-level suppliers when necessary to overcome failure, consistent poor quality, or other problems. Audits shall be planned for and performed using established and maintained procedures (see ECSS-Q-ST-10, clause 5.2.3).} &
        &
        {SPAP 6.8} &
        \\ \hline

        SPA-05 &
        {\textbf{Software problems} \newline
        Procedures for the logging, analysis and correction of all software problems encountered during software verification and validation shall be defined and implemented.  \newline
        Software problem reports during verification and validation shall contain the following information:  \newline
        • identification of the software item;  \newline
        • description of the problem;  \newline
        • recommended solution;  \newline
        • final disposition;  \newline
        • modifications implemented (e.g. documents, code, and tools);  \newline
        • tests re-executed.} &
        &
        {SPAP 6.8, \newline SPAMR 9} &
        \\ \hline

        SPA-06 &
        {\textbf{Methods and tools} \newline
        Methods and tools to be used for all the activities of the software development cycle (including requirements analysis, software specification, modelling, design, coding, validation, testing, configuration management, verification and product assurance) shall be identified by the supplier and agreed by the customer. \newline
        The correct use of methods and tools shall be verified and reported.} &
        &
        {SPAP 5.8, \newline SPAMR 5} &
        \\ \hline

        SPA-07 &
        {\textbf{Documentation of processes} \newline
        The following activities shall be covered by project planning documentation: \newline
        • development; \newline
        • specification, design and customer documents to be produced; \newline
        • configuration and documentation management; \newline
        • verification, testing and validation activities; \newline
        • maintenance. \newline
        Project plans, procedures and standards shall be finalized before the start of the related activities.} &
        &
        {SPAP 6.2} &
        \\ \hline

        SPA-08 &
        {\textbf{Software dependability} \newline
        The software shall be specified and designed to prevent failure propagation. \newline
        Error handling shall be implemented to ensure that, in case of fault, the software shall support graceful degradation.} &
        &
        {SPAP 6.3} &
        \\ \hline

        SPA-09 &
        {\textbf{Software configuration management} \newline
        Besides the configuration requirements applicable to the project, the following specific requirements apply to software configuration management: \newline
        • the software configuration management system shall allow any reference version to be re-generated from backups; \newline
        • the supplier shall ensure that all authorized changes are implemented in accordance with the software configuration management planning and procedures; \newline
        • software configuration management planning and procedures shall address software branching and merging.} &
        &
        {SPAP 6.4} &
        \\ \hline

        SPA-10 &
        {\textbf{Process metrics} \newline
        The following basic process metrics shall be used: \newline
        • duration: how phases and tasks are being completed versus the planned schedule; \newline
        • effort: how much effort is consumed by the various phases and tasks compared to the plan. \newline
        The following basic process metric shall be used and reported to the customer: \newline
        • Number of problems detected during validation testing} &
        &
        {SPAP 6.5, \newline SPAMR 7.2} &
        \\ \hline

        SPA-11 &
        {\textbf{Verification} \newline
        The software assurance personnel shall ensure that: \newline
        • activities for the verification of the quality requirements are specified; \newline
        • verification activities are performed according to the plan; \newline
        • verification is carried out by suitably independent personnel; \newline
        • the outputs of each development activity are verified for conformance against pre-defined criteria, and only outputs which have been subjected to planned verifications are used as inputs for subsequent activities; \newline
        • the completion of actions related to software problem reports generated during verification is verified and recorded; \newline
        • traceability matrices are verified at each milestone.} &
        &
        {SPAP 6.7.5, \newline SPAMR 4} &
        \\ \hline

        SPA-12 &
        {\textbf{Reuse of existing software} \newline
        The existing software re-used within the project shall be identified and agreed with ESA.} &
        C &
        {SPAP 6.6} &
        \\ \hline

        SPA-13 &
        {\textbf{Software requirements} \newline
        Besides functional requirements, the following classes of non-functional requirements shall be specified: \newline
        • performance, \newline
        • reliability, \newline
        • robustness, \newline
        • quality, \newline
        • maintainability, \newline
        • configuration management, \newline
        • security, \newline
        • privacy, \newline
        • metrication, and \newline
        • verification and validation. \newline
        \newline
        The Supplier shall agree on the following principles and rules: \newline
        • assignment of persons (on both sides) responsible for establishing the software requirements; \newline
        • methods for agreeing on requirements and approving changes; \newline
        • efforts to prevent misunderstandings such as definition of terms, explanations of background of requirements; \newline
        • recording and reviewing discussion results on both sides.} &
        C &
        {SRS, \newline SPAP 6.7} &
        \\ \hline

        SPA-14 &
        {\textbf{Software architectural design and design of software items} \newline
        It shall be ensured that: \newline
        • mandatory and advisory design standards are defined and applied, adherence to design standards is verified; \newline
        • specific rules on design and code are defined to ensure that the specified level of numerical accuracy is obtained; \newline
        • means, criteria and tools to ensure that the complexity and modularity of the design meet the quality requirements are defined and applied; \newline
        • the design documentation contains the appropriate level of information for maintenance activities.} &
        &
        {SPAP 6.7, \newline SPAMR 6.1} &
        \\ \hline

        SPA-15 &
        {\textbf{Coding} \newline
        It shall be ensured that: \newline
        • coding standards (including consistent naming conventions and adequate commentary rules), consistent with the product quality requirements, are specified and observed; \newline
        • the tools to be used in implementing and checking conformance with coding standards are identified before coding activities start; \newline
        • the code evaluation against the coding standards and applicable quality requirements is performed in parallel with the coding process, in order to provide feedback to the software programmers; \newline
        • synthesis of the code analysis results and corrective actions implemented are described in the software product assurance reports; \newline
        • the software code is put under configuration control immediately after successful unit testing.} &
        &
        {SPAP 6.7, \newline SPAMR 6.2} &
        \\ \hline

        SPA-16 &
        {\textbf{Testing and validation} \newline
        It shall be ensured that: \newline
        • the test procedures and data are adequate, feasible and traceable and that they satisfy the requirements; \newline
        • test readiness reviews shall be held before the commencement of test activities; \newline
        • software problem detected during testing are properly documented and reported to those concerned; \newline
        • the completion of actions related to software problem reports generated during testing and validation is verified and recorded; \newline
        • provisions are made to allow witnessing of tests by the customers; \newline
        • tests are conducted in accordance with approved test procedures and data; \newline
        • the configuration under test is correct; \newline
        • the tests are properly documented; \newline
        • the test reports are up to date and valid; \newline
        • that tests are repeatable, by verifying the storage and recording of tested software, support software, test environment, supporting documents and problems found; \newline
        • areas affected by any modification are identified and re-tested (regression testing); \newline
        • test documentation is up to date and organized to facilitate its reuse for maintenance; \newline
        • before offering the product for delivery and customer acceptance, the software is validated as a complete product, under conditions similar to the application environment as specified in the software requirements.} &
        &
        {SPAP 6.7} &
        \\ \hline

        SPA-17 &
        {\textbf{Software delivery and acceptance} \newline
        The roles, responsibilities and obligations of the supplier and customer during installation shall be established. \newline
        The Supplier shall ensure that: \newline
        • the delivered software complies with the contractual requirements (including any specified content of the software acceptance data package); \newline
        • the source and object code supplied correspond to each other; \newline
        • all agreed changes are implemented; \newline
        • all software problems are either resolved or declared. \newline
         \newline
        The customer shall: \newline
        • establish an acceptance test plan, specifying the intended acceptance tests; \newline
        • ensure that the acceptance tests are performed in accordance with the approved acceptance test plan; \newline
        • verify that the executable code was regenerated from configuration managed source code components and installed in accordance with predefined procedures on the target environment; \newline
        • ensure that any discovered problems are documented; \newline
        • certify conformance to the procedures and state the conclusion concerning the acceptance result for the software product under test (accepted, conditionally accepted, rejected).} &
        &
        {SDP, \newline SPAP 6.7} &
        \\ \hline

        SPA-18 &
        {\textbf{Product quality objectives and metrication} \newline
        A quality model shall be defined, based on the following quality characteristics: \newline
        • functionality; \newline
        • reliability; \newline
        • maintainability; \newline
        • usability. \newline
         \newline
        The quality model shall be used to specify software quality requirements in quantitative terms or constraints. A metrication programme shall be defined, based on the quality model, specifying: \newline
        • the metrics to be collected and stored; \newline
        • the means to collect metrics (measurements); \newline
        • the target values, with reference to the product quality requirements; \newline
        • how the results of the analyses performed on the collected metrics are fed back to the development team and used to identify corrective actions; \newline
        • the schedule of metrics collection, storing, analysis and reporting, with reference to the whole software life cycle. \newline
         \newline
        Numerical accuracy shall be estimated and verified. \newline
        The results of metrics collection and analysis shall be included in the software product assurance reports.} &
        &
        {SPAP 5.5, 7, \newline SPAMR 7.1} &
        \\ \hline

        SPA-19 &
        {\textbf{Quality of software requirements} \newline
        The software requirements shall be: \newline
        • correct; \newline
        • unambiguous; \newline
        • complete; \newline
        • consistent; \newline
        • verifiable; \newline
        • traceable. \newline
        For each requirement, the method for verification and validation shall be specified.} &
        &
        {SRS, ?} &
        \\ \hline

        SPA-20 &
        {\textbf{Quality of design and related documentation} \newline
        The software design shall: \newline
        • meet the non-functional requirements; \newline
        • engineered in a way to facilitate testing; \newline
        • be subjected to no dependency on the operating system and the hardware, in order to aid portability} &
        &
        {SDD, \newline SPAP 6.7.2} &
        \\ \hline

        SPA-21 &
        {\textbf{Quality of test and validation documentation} \newline
        The test documentation shall: \newline
        • cover the test environment, tools and test software; \newline
        • specify the criteria for completion of each test and any contingency steps; \newline
        • specify test procedures, data and expected results; \newline
        • identify hardware and software configuration for testing.} &
        &
        {SIVVP ?} &
        \\ \hline
        
    \end{xltabular}
}

%{
%	\renewcommand{\arraystretch}{1.4} % introduce padding on top and bottom of cells
%	\tiny
%	\fontfamily{phv}\selectfont
%	
%	\begin{xltabular}{\textwidth}{|l|l|S|L|l|S|}
%	
%	% Header on all pages
%	\hline
%	ID & Clause & Title & Applicable Requirements & Compliance & Document \\ \hline\hline
%	\endhead
%	
%	% CSV file content
%	\csvreader[head to column names,
%    	       separator=semicolon,
%    	       late after line = \\\hline]
%    	       {SPAP_Traceability.csv}{}
%    	       {\csvcoli & \csvcolii & \csvcoliii &
%    	       \csvcoliv & \csvcolv & \csvcolvi}
%	\end{xltabular}
%}